\chapter{Intersex Variation}

\section*{Definition and Overview}
Intersex is an umbrella term for people born with variations in sex characteristics: chromosomes, hormones, genitals, or reproductive organs that do not fit typical definitions of male or female. Intersex variation is sometimes called differences of sex development (DSD). It is a natural form of human diversity, not a disease, and most intersex people are healthy. Intersex variations include conditions such as Klinefelter syndrome, Turner syndrome, congenital adrenal hyperplasia, and androgen insensitivity syndrome. Care should respect the person's autonomy, avoid unnecessary and irreversible interventions on infants who cannot consent, and provide accurate information, psychological support, and medical care where it is genuinely needed.

\section*{Epidemiology}
Intersex variations occur in around one in 1500 to 1 in 2000 live births, depending on which variations are counted, making them roughly as common as red hair. Some variations, such as congenital adrenal hyperplasia, are detected in newborn screening in Australia. Intersex people are found in all cultures and communities. Historically, many intersex infants underwent surgery and hormone treatment without consent to make their bodies conform to typical male or female anatomy, but this practice has been strongly criticised, and Australian guidelines now emphasise delaying irreversible procedures until the person can decide for themselves.

\section*{Aetiology and Pathophysiology}
Sex development is controlled by chromosomes, genes, and hormones acting in sequence. Variations arise from differences at any of these levels: sex chromosomes that are neither XX nor XY, such as XXY, XO, or mosaic patterns; gene variants affecting hormone production or action, such as androgen insensitivity or 5-alpha reductase deficiency; or differences in the development of the gonads and genitals. The result can be a body that develops along a pathway different from typical male or female development. Most intersex variations cause no health problems, but some are associated with medical needs such as hormone replacement, fertility considerations, or specific health risks that require monitoring.

\section*{Clinical Features}
\begin{itemize}
    \item Genitals that appear atypical at birth, though many intersex variations are not visible
    \item Variations in chromosomes, such as XXY or XO, found on testing
    \item Differences in puberty: it may not start, may start late, or may follow an unexpected pattern
    \item Differences in fertility and the need for assisted reproduction or hormone therapy
    \item Some conditions have associated health effects, such as bone, heart, or kidney considerations
    \item Congenital adrenal hyperplasia may cause salt loss and require urgent treatment in infancy
    \item Many intersex people have no symptoms and live healthy lives
    \item Identity and wellbeing vary; most intersex people identify as male or female, and some as non-binary
\end{itemize}

\section*{Differential Diagnosis}
Assessment of an intersex variation aims to identify the specific condition and its medical implications.
\begin{itemize}
    \item \textbf{Chromosomal variations:} Klinefelter, Turner, and other chromosome patterns
    \item \textbf{Androgen insensitivity syndrome:} XY chromosomes with a body that does not respond to testosterone
    \item \textbf{Congenital adrenal hyperplasia:} an enzyme problem affecting hormone production
    \item \textbf{5-alpha reductase deficiency and other enzyme conditions:} affecting hormone conversion
    \item \textbf{Ovotesticular and gonadal variations:} differences in the gonads themselves
    \item \textbf{Other causes of atypical genitals:} rare syndromes and maternal hormone effects
\end{itemize}

\section*{When to Seek Medical Care}
Families of a newborn with an intersex variation should be supported by a specialist team that includes a paediatrician, endocrinologist, urologist, geneticist, and psychologist, with accurate, balanced information and no pressure for rushed decisions. Adults with intersex variations should have their health needs reviewed, including bone, hormone, and reproductive health. Anyone experiencing distress about their sex characteristics should have access to counselling.

\textbf{Call 000 (or local emergency services) for:}
\begin{itemize}
    \item A newborn with signs of salt-wasting crisis: poor feeding, vomiting, lethargy, or dehydration, which may indicate congenital adrenal hyperplasia
    \item Severe pain or acute illness in an infant with a known intersex condition
    \item Any medical emergency in a person with an intersex variation and associated health conditions
\end{itemize}

\section*{Diagnosis and Investigations}
\begin{itemize}
    \item \textbf{Karyotype:} chromosome testing to identify sex chromosomes
    \item \textbf{Hormone tests:} measuring testosterone, oestrogen, adrenal, and other hormone levels
    \item \textbf{Genetic testing:} for specific gene variants causing intersex conditions
    \item \textbf{Imaging:} ultrasound to assess internal reproductive organs
    \item \textbf{Newborn screening:} detects some forms of congenital adrenal hyperplasia
    \item \textbf{Specialist team assessment:} endocrinology, genetics, urology, and psychology
    \item \textbf{Health monitoring:} bone density, heart, kidney, and other checks according to the condition
\end{itemize}

\section*{Management}
\begin{itemize}
    \item \textbf{Accurate information and counselling:} provided to the person and family without bias
    \item \textbf{Delaying irreversible interventions:} avoiding surgery on infants and young children until they can consent
    \item \textbf{Hormone therapy:} replacement hormones where needed for puberty, bone health, or wellbeing
    \item \textbf{Treatment of associated conditions:} such as salt-wasting in congenital adrenal hyperplasia
    \item \textbf{Fertility support:} advice and assisted reproduction where possible
    \item \textbf{Psychological support:} for identity, body image, and family adjustment
    \item \textbf{Peer support:} connection with intersex community organisations
    \item \textbf{Medical follow-up:} regular review of condition-specific health needs
\end{itemize}

\section*{Complications}
\begin{itemize}
    \item Harm from non-consensual early surgery, including loss of sensation, scarring, and psychological trauma
    \item Stigma, discrimination, and misunderstanding
    \item Infertility and its psychological impact
    \item Condition-specific health risks, such as osteoporosis, heart problems, or salt-wasting crisis
    \item Anxiety, depression, and identity distress when support is inadequate
    \item Barriers to affirming, knowledgeable healthcare
\end{itemize}

\section*{Prognosis and Outcomes}
With respectful, informed care, intersex people live healthy and fulfilling lives. The shift away from non-consensual infant surgery, supported by Australian and international guidelines, protects bodily autonomy and improves long-term psychological outcomes. Medical care is available for the specific health needs that some intersex variations carry, and most are manageable. Peer support and community connection are powerful protective factors. The most important elements are accurate information, psychological support, and healthcare that treats intersex as diversity rather than defect.

\section*{Prevention and Self-Care}
\begin{itemize}
    \item Seek care from health professionals with experience and knowledge of intersex variation
    \item Ask for a specialist multidisciplinary team when a variation is found
    \item Do not rush decisions about surgery or treatment for an infant
    \item Obtain a second opinion and accurate information before any irreversible procedure
    \item Connect with intersex peer support organisations
    \item Attend recommended health monitoring for your condition
    \item Seek counselling for identity, body image, or relationship concerns
    \item Speak up if healthcare is not respectful or knowledgeable
    \item Advocate for consent-based care in health services
\end{itemize}

\section*{Sources and Further Reading}
The following reputable sources provide further information for healthcare students and patients seeking reliable, current advice about intersex variation.
\begin{itemize}
    \item Intersex Human Rights Australia. \textit{Information and advocacy.} https://ihra.org.au/
    \item Department of Health and Aged Care. \textit{Guidelines on the health of intersex people.} https://www.health.gov.au/
    \item Androgen Insensitivity Syndrome Support Group Australia. \textit{Peer support and information.} https://aissga.org.au/
    \item Healthdirect Australia. \textit{Intersex variations.} https://www.healthdirect.gov.au/
    \item Royal Children's Hospital Melbourne. \textit{Differences of sex development.} https://www.rch.org.au/
    \item National Organization for Rare Disorders. \textit{Intersex and DSD conditions.} https://rarediseases.org/
\end{itemize}
